\documentclass[11pt]{article}
\usepackage{mathpartir}
\usepackage{pifont}
\usepackage{multicol}
\usepackage{listings}
\usepackage{pgf}
\usepackage{tikz}
\usepackage{alltt}
\usepackage{hyperref}
\usepackage{url}
\usepackage{amsmath}
\usepackage{amssymb}
\usepackage{dafny}
\usetikzlibrary{arrows,automata,shapes,positioning}
\tikzstyle{block} = [rectangle, draw, fill=blue!20, 
    text width=5em, text centered, rounded corners, minimum height=2em]
\tikzstyle{bt} = [rectangle, draw, fill=blue!20, 
    text width=1em, text centered, rounded corners, minimum height=2em]
\newcommand{\xmark}{\ding{55}}

\newtheorem{defn}{Definition}
\newtheorem{crit}{Criterion}
\newcommand{\true}{\mbox{\sf true}}
\newcommand{\false}{\mbox{\sf false}}

\newcommand*\circled[1]{\tikz[baseline=(char.base)]{
            \node[shape=circle,draw,inner sep=2pt] (char) {#1};}}


\newcommand{\handout}[5]{
  \noindent
  \begin{center}
  \framebox{
    \vbox{
      \hbox to 5.78in { {\bf Software Testing, Quality Assurance and Maintenance } \hfill #2 }
      \vspace{4mm}
      \hbox to 5.78in { {\Large \hfill #5  \hfill} }
      \vspace{2mm}
      \hbox to 5.78in { {\em #3 \hfill #4} }
    }
  }
  \end{center}
  \vspace*{4mm}
}

\newcommand{\lecture}[4]{\handout{#1}{#2}{#3}{#4}{Lecture #1}}
\topmargin 0pt
\advance \topmargin by -\headheight
\advance \topmargin by -\headsep
\textheight 8.9in
\oddsidemargin 0pt
\evensidemargin \oddsidemargin
\marginparwidth 0.5in
\textwidth 6.5in

\parindent 0in
\parskip 1.5ex
%\renewcommand{\baselinestretch}{1.25}

\usepackage{enumitem}

\newtheorem{prop}{Proposition}
\newtheorem{lemma}{Lemma}
\usepackage{ebproof}
\newcommand{\qedsymbol}{\rule{1ex}{1ex}}
\newcommand{\sem}[3]{\langle #1, #2 \rangle \Downarrow #3}

\lstset{ %
language=Java,
basicstyle=\ttfamily,commentstyle=\scriptsize\itshape,showstringspaces=false,breaklines=true,numbers=left}

%\usepackage{fontspec}
%\setmonofont{Cousine}[Scale=MatchLowercase]

\begin{document}

\lecture{19 --- March 16, 2026}{Winter 2026}{Patrick Lam}{version 2}

\section*{Loops and Loop Invariants}
The usual more efficient way to compute Fibonacci numbers is using a loop. You've seen loops in the context of Hoare logic, and Dafny is similar, except that it does a lot of
proving for you. But it doesn't supply invariants.

Here's a loop, with an invariant.
\begin{lstlisting}[language=dafny]
method FirstLoop(n:nat)
{
  var i := 0;
  while i < n
    invariant 0 <= i
  {
    i := i + 1;
  }
  assert i == n;
}
\end{lstlisting}
We can see that $\mathsf{i == n}$ at the end of the loop, but does Dafny know that? We can ask it, using
\textsf{assert}. (It doesn't). We need to strengthen the loop invariant. If we try $0 \leq i < n$, then
Dafny complains that the invariant might not hold on entry and it might not be maintained by the loop body.
Why is that?

Well, in any case, if we use the invariant $0 \leq i \leq n$, then Dafny is satisfied with everything.

\paragraph{Exercise 6.} Change the loop invariant to $0 \leq i \leq n + 2$. Does the loop still verify?
Does the assertion $i == n$ after the loop still verify?

\paragraph{Exercise 7.} With the original loop invariant, change the loop guard from $i < n$ to
$i \neq n$. Do the loop and assertion after the loop still verify? Why or why not?

Here is another loop, which implements multiplication using addition.
\begin{lstlisting}[language=dafny]
method AddByInc(n:nat, m:nat) returns (r:nat)
  ensures r == n + m
{
  r := n;
  var i := 0;
  while (i < m)
  {
    i := i + 1;
    r := r + 1;
  }
}
\end{lstlisting}

Turns out that we need to add two invariants for the while:
\begin{lstlisting}[language=dafny]
    invariant 0 <= i <= m;
    invariant r == n + i;
\end{lstlisting}
These are similar invariants to what we've seen in Hoare logic, but it's easier to experiment with them here and see exactly what we need and when Dafny complains.
The second invariant really is an invariant. We can express it as $n = r - i$. The loop doesn't modify $n$. So, as $i$ goes up, $r$ goes down by the same amount.

We can also check some assertions at the end to see what Dafny actually knows:
\begin{lstlisting}[language=dafny]
  assert (r == n + i && i == m);
  assert (r == n + m);
\end{lstlisting}

Here's a method that multiplies two numbers by addition.
\begin{lstlisting}[language=dafny]
method Product(m:nat, n:nat) returns (res:nat)
  ensures res == m*n
{
  var m1 := m;
  res := 0;
  while m1 != 0
  {
    var n1 := n;
    while n1 != 0
    {
      res := res + 1;
      n1 := n1 - 1;
    }
    m1 := m1 - 1;
  }
}
\end{lstlisting}
We need invariants for both \textsf{while} loops. We can start with the outer loop.
We could write an invariant about the range of \textsf{m1} but it turns out that this is not
mandatory. The invariant that actually is mandatory is about the value of \textsf{res}.
It looks like magic, but you can examine the postcondition and the fact that \textsf{m1}
is the induction variable here, and conclude that your invariant should be:
\begin{lstlisting}[language=dafny]
  invariant res == (m-m1) * n
\end{lstlisting}
When the loop starts, it has gone around the loop zero times, so
$\mathsf{m} == \mathsf{m1}$ and their difference is 0, as is
$\mathsf{res}$.  In terms of maintaining the invariant, each time around
the outer loop, the program is adding one \textsf{n} to $\mathsf{res}$
and subtracting 1 from \textsf{m1}, and the invariant does that.

As for the second invariant, we are now doing the adding of \textsf{n} to \textsf{res} one at a time.
Through this loop, \textsf{m} and \textsf{m1} are both constant, and we are changing \textsf{res}
from $(\mathsf{m}-\mathsf{m1}) \times n$ to $(\mathsf{m}-\mathsf{m1}) + 1 \times n$. We have that \textsf{n} stays constant
and \textsf{n1} goes down. Every time that \textsf{n1} goes down by 1, \textsf{res} goes up by 1,
so putting $-\mathsf{n1}$ is a good bet.
The invariant is thus
\begin{lstlisting}
  invariant res == (m-m1)*n + (n-n1)
\end{lstlisting}
and Dafny will verify this.

Let's implement the Fibonacci method.
\begin{lstlisting}[language=dafny]
method ComputeFib(n: nat) returns (b: nat)
  ensures b == fib(n)
{
  var i := 1;
  var a := 0;
  b := 1;
  while i < n
  {
    a, b := b, a + b;
    i := i + 1;
  }
}
\end{lstlisting}
(Note that Dafny allows assignment to \textsf{a} and \textsf{b} simultaneously).

We have postcondition $b = \mathsf{fib}(n)$, and we also know
that $i == n$, so it must be that $b == \mathsf{fib}(i)$ which we can have as a loop invariant. This
is promising as part of the loop invariant: it talks about the loop counter. Hence, we have what we had before,
plus this.
\begin{lstlisting}[language=dafny]
  invariant 0 <= i <= n
  invariant b == fib(i)
\end{lstlisting}
We also know that $a$ is the previous Fibonacci number. So,
\begin{lstlisting}[language=dafny]
  invariant a == fib(i - 1)
\end{lstlisting}
Finally, since Dafny \textsf{nat} starts at 0, we include:
\begin{lstlisting}[language=dafny]
  if n == 0 { return 0; }
\end{lstlisting}
All together:
\begin{lstlisting}[language=dafny]
function fib(n:nat): nat
{
  if n == 0 then 0
  else if n == 1 then 1
  else fib(n-1) + fib(n-2)
}

method ComputeFib(n:nat) returns (b:nat)
  ensures b == fib(n)
{
  if n == 0 { return 0; }
  var i := 1;
  var a := 0;
  b := 1;
  while i < n
    invariant 0 < i <= n
    invariant a == fib(i - 1)
    invariant b == fib(i)
  {
    a, b := b, a + b;
    i := i + 1;
  }
}
\end{lstlisting}

% stopped about here after 90minutes

\paragraph{Exercise 8.} It is possible to simplify \textsf{ComputeFib}.
Write a simpler program by not introducing \textsf{a} as the Fibonacci number
that precedes \textsf{b}, but instead introducing \textsf{c} as the variable that
succeeds \textsf{b}. Verify that your program is correct according to the mathematical
definition of Fibonacci.

(I'm omitting Exercise 9 from the tutorial because it's not really that useful.)

As you know, an issue with just verifying partial correctness is that it doesn't tell you
anything at all if the program doesn't terminate. When you run the program, you can see it looking like it's going to not
terminate (though how do you know for sure?), but if you're just proving it, you might write non-terminating code.

\section*{Slow exponentiation}
Here's another method with a loop. This method computes $2^N$.
Let's start with a general-purpose function that computes $x^n$
where $x \in \mathbb{Z}$ and $n \in \mathbb{N} \cup \{ 0 \}$.

\begin{lstlisting}[language=dafny]
function exp(x:int, n:nat):int
{
  if n == 0 then 1 else if n == 1 then x else x * exp(x, n-1)
}
\end{lstlisting}

And some code to repeatedly double $p$.
\begin{lstlisting}[language=dafny]
method PowerOfTwo(N:nat) returns (p:int)
{
  var r : int := 0;
  var i : nat := 0;
  p := 1;
  while (i < N)
  {
    p := 2 * p;
    i := i + 1;
  }
}
\end{lstlisting}

Looking at the method name and the implementation, we can deduce
this \textsf{ensures} clause:
\begin{lstlisting}[language=dafny]
  ensures p == exp(2, N)
\end{lstlisting}
and the usual invariant on the induction variable:
\begin{lstlisting}[language=dafny]
  invariant i <= N
\end{lstlisting}
as well as one involving $p$:
\begin{lstlisting}[language=dafny]
  invariant p == exp(2, i)
\end{lstlisting}

\end{document}
